\section{张量积}

\subsection{线性映射的张量积}

\begin{definition}
    \

    给出 $F$ 向量空间 $V,W,V',W'$. 取线性映射 $\varphi:V \to V', \psi:W \to W'$. 定义映射 $\varphi \odot  \psi: V \otimes W \to V' \otimes W'$:
    \vspace{-1em}
    \begin{equation*}
        (\varphi \odot \psi)(v \otimes w) = \varphi(v) \otimes \psi(w).
    \end{equation*}
\end{definition}

显然 $(\varphi \odot \psi)(v,w) = \varphi(v) \otimes \psi(w)$ 是双线性映射, 由张量的泛性质保证出 $\varphi \odot \psi$ 的线性性. 这诱导出另一个映射
\begin{equation*}
   \begin{matrix}
        g:\text{Hom}_F(V,V') \times \text{Hom}_F(W,W') \to \text{Hom}_F(V \otimes W, V' \otimes W'),\\ 
        \varphi,\psi \mapsto \varphi \odot \psi.
    \end{matrix} 
\end{equation*}
$g$ 也是双线性映射.

对 $\text{Hom}_F(V,V')$ 和 $\text{Hom}_F(W,W')$ 作张量积将诱导出如下的交换图

\begin{center}
    \begin{tikzcd}
        \text{Hom}_F(V,V') \times \text{Hom}_F(W,W') \arrow[rd, "g"] \arrow[r] 
        & \text{Hom}_F(V,V') \otimes \text{Hom}_F(W,W') \arrow[d, dashed, "h"] \\
        & \text{Hom}_F(V \otimes W, V' \otimes W')
    \end{tikzcd}
\end{center}
% \vspace{0.6em}
\begin{theorem}
  \

  存在唯一线性映射
  \begin{equation*}
    \begin{matrix}
        h:\text{Hom}_F(V,V') \otimes \text{Hom}_F(W,W') \to \text{Hom}_F(V \otimes W, V' \otimes W'),\\ 
        \varphi \otimes \psi \mapsto \varphi \odot \psi.
    \end{matrix} 
\end{equation*}
且 $h$ 是单射. 如果所有向量空间都是有限维, 则 $h$ 是同构.
\end{theorem}

\hspace{-2em}\textbf{证} \hspace{0.5em} 只需证明 $\ker{h} = \{0\}$. 

采用反证法, 假设 $\eta \in \text{Hom}_F(V,V') \otimes \text{Hom}_F(W,W')$ 且 $\eta \neq 0$, 但 $h(\eta) = 0$. 取 $\text{Hom}_F(W,W')$ 中的一组线性无关基底 $\{\psi_1, \psi_2, \cdots, \psi_n\}$, 以及 $\{\varphi_i\} \subset \text{Hom}_F(V,V')$, 使得
\begin{equation*}
    \eta = \sum_{i=1}^n \varphi_i \otimes \psi_i.
\end{equation*}

对于任意 $v \in V, w \in W$, 由 $h(\eta) = 0$ 可得
\begin{equation*}
    0 = h(\eta)(v \otimes w) =h \left( \sum_{i=1}^n \varphi_i \otimes \psi_i \right)(v \otimes w)= \sum_{i=1}^n \varphi_i(v) \otimes \psi_i(w).
\end{equation*}

由于 $\eta \ne 0$, 存在一 $\varphi_i \ne 0$. 所以有 $\varphi_i(v) \ne 0$ 对于某个 $v \in V$, 我们假定 $\varphi_1(v')\ne 0$. 从 $\{\varphi_i\}_{1 \ge i \ge n}$ 中选出最大线性无关子集 $\{\varphi_i\}_{1 \ge i \ge m}$ (我们假定重排成恰好该情况). 那么
\begin{align*}
    0 &= \sum_{i=1}^n \varphi_i(v') \otimes \psi_i(w) \\
    &= \sum_{i=1}^m \varphi_i(v') \otimes \psi_i(w) + \sum_{k=m+1}^n \left (\sum_{j=1}^m a_{k,j}\varphi_j(v') \right) \otimes \psi_k(w) \\
    &= \sum_{i=1}^m \varphi_i(v') \otimes w_i'.
\end{align*}
其中 $w_i' = \psi_i(w) + \sum_{k=m+1}^n a_{k,i}\psi_k(w)$. 由于 $\{\varphi_i(v')\}_{1 \ge i \ge m}$ 线性无关, 故 $w_i' = 0$, 即
\begin{equation*}
    \psi_i(w) + \sum_{k=m+1}^n a_{k,i}\psi_k(w) = 0
\end{equation*}
这与集合 $\{\psi_i\}$ 线性无关性矛盾, 因此 $h$ 是单射. 由于空间都是有限维, 只需证明两侧维数相等即得出同构, 在此省略. \hfill\(\qed\)

此后便不再区分 $\varphi \odot \psi$ 和 $\varphi \otimes \psi$, 例如 $\varphi \otimes \psi(v \otimes w) = \varphi(v) \otimes \psi(w)$.

下面的例子说明了两个线性映射的张量积对应于相应矩阵的 Kronecker 积.

\vspace{0.6em}

\begin{example}
  \

  取 $n \times n$ 矩阵 $A_{ij}$ 和 $m \times m$ 矩阵 $B_{ij}$. 给出与之相对应的两个线性映射
  \begin{equation*}
    \varphi_A: F^n \to F^n, \varphi_B: F^m \to F^m, 
  \end{equation*}
  记作 $[\varphi_A]=A, [\varphi_B]=B$ (矩阵对应于 $F^n$ 的标准基 $\{e_i\}$). 

  如果在 $F^n \otimes F^m$ 取基 (字典序)
  \begin{equation*}
    e_1 \otimes f_1, \cdots, e_1 \otimes f_m, \cdots, e_2 \otimes f_1, \cdots, e_2 \otimes f_m, \cdots, e_n \otimes f_1, \cdots, e_n \otimes f_m,
  \end{equation*}
  那么 $[\varphi_A \otimes \varphi_B] = A \otimes B$, 后者是矩阵的 Kronecker 积.
\end{example}

\hspace{-2em}\textbf{证} \hspace{0.5em} 将线性映射作用在地 $k$ 个基向量 $e_p \otimes f_q$ 上
\begin{align*}
    (\varphi_A \otimes \varphi_B)(e_p \otimes f_q) &= \varphi_A(e_p) \otimes \varphi_B(f_q) \\
    &= \left( \sum_{i=1}^n A_{ip} e_i \right) \otimes \left( \sum_{j=1}^m B_{jq} f_j \right) \\
    &= \sum_{i=1}^n\sum_{j=1}^m \left(A_{ip}B_{jq} e_i \otimes f_j \right).
\end{align*}
这给出了矩阵第 $(p-1)m+q$ 列的元素.

也就是说, 矩阵的第 $(i-1)m+j$ 行, 第 $(p-1)m+q$ 列的元素恰好是 $A_{ip}B_{jq}$. 这与 Kronecker 积的定义一致. \hfill\(\qed\)


\newpage